\chapter*{Мобайл AR аялал жуулчлалын тоглоомын систем дизайны онолын тойм}
\begin{sloppy}

\section*{Гүйцэтгэх хураангуй}
Мобайл AR аялал жуулчлалын тоглоомын үндсэн инженерчлэлийн асуудал нь локал AR поз ба глобал геопозицыг нэгэн зэрэг зөв удирдах асуудал юм. GPS/GNSS нь хэрэглэгчийн глобал байршлыг өгдөг боловч ухаалаг утас нээлттэй тэнгэрийн дор ч түгээмэлдээ ойролцоогоор 4.9 м радиустай алдаатай байдаг. Иймээс сайн архитектур нь географ trigger давхаргыг визуал/инерцийн бэхэлгээ давхаргаас тусгаарлах ёстой.

Платформ тодорхойгүй үед хамгийн зөв чиглэл нь capability-based, degradation-friendly архитектур юм. Клиент талд нэг ``AR session orchestration'' давхарга байрлуулж, мэдрэгчийн бэлэн байдал, сүлжээний чанар зэргийг хэмжин тухайн мөчийн тоглоомын горимыг сонгодог байх нь зөв.

\section*{Суурь архитектурын зарчим}
Логикийг дөрвөн том хавтгайд салгах нь оновчтой:
\begin{itemize}
    \item \textbf{Perception plane:} Камер ба мэдрэгчээс local pose үүсгэнэ.
    \item \textbf{Geospatial plane:} Map/route/geofence/POI логикийг шийднэ.
    \item \textbf{Content plane:} Localized narrative, quest, media asset-уудыг удирдана.
    \item \textbf{Service plane:} User state, realtime sync, analytics-ийг хариуцна.
\end{itemize}

\section*{Клиент талын мэдрэгч ба AR Perception Stack}
Гадаах AR аяллын тоглоомд нэг ч мэдрэгч хангалттай биш. Архитектурын хувьд дараах мэдрэгчүүдийн хослолыг ашигладаг:
\begin{itemize}
    \item \textbf{GPS/GNSS:} Глобал байршил, бүсийн trigger , POI prefetch-д зориулагдсан.
    \item \textbf{IMU (Accelerometer + Gyroscope):} Өндөр давтамжтай хөдөлгөөн, pose smoothing хийхэд ашиглагдана.
    \item \textbf{Magnetometer:} Magnetic north-той холбосон heading буюу чиглэл тогтооход тусална.
    \item \textbf{Barometer:} Өндрийн ялгаа, шаталсан орчинд trigger өгөхөд ашиглагдана.
    \item \textbf{Камер + SLAM/VIO:} Локал pose, гадаргуу илрүүлэлт, stable anchors үүсгэх үндсэн хэрэгсэл.
\end{itemize}

\section*{Геолокаци ба Орон зайн логик}
Глобал ба локал хоёр координатын ертөнцийг салгаж ойлгох нь чухал. Глобал ертөнц нь lat/lon/alt-аар бодогддог бол локал ертөнц нь camera pose, anchor-оор бодогдоно. 2024 оны судалгаагаар LBAR системүүд хэт хатуу маршрутаар хэрэглэгчийг чиглүүлбэл орон зайн танин мэдэхүйн эрх чөлөөг хязгаарлах талтай байдаг тул detour-friendly буюу уян хатан маршрут төлөвлөх нь зөв.

\section*{Контент хүргэлт ба Серверийн архитектур}
Контентыг metadata (JSON/Proto) ба binary asset (GLB, texture, audio) гэсэн хоёр өөр ангиар удирдах нь оновчтой. Том медиа файлуудыг CDN-ээр дамжуулан хүргэж, хэрэглэгчийн төлөвийг (state) authoritative backend-д хадгална.
\begin{itemize}
    \item \textbf{Offline-first стратеги:} Зөвхөн газрын зураг бус, ойролцоох POI, localized text, 3D bundles зэргийг local cache-д хадгалах.
    \item \textbf{Headless CMS:} Presentation-ийг backend-ээс салгаж, олон хэлний (localization)Fallback логикийг хэрэгжүүлэх.
\end{itemize}

\section*{Нууцлал ба Талбайн QA}
Байршил, камер нь мэдрэмтгий өгөгдөл тул Least Privilege зарчмыг баримталж, зөвшөөрлийг progressive disclosure хэлбэрээр (яг хэрэгцээ гарсан үед нь) авах ёстой. Туршилтын хувьд ARCore Recording \& Playback API ашиглан бодит орчны session-ийг бичиж, deterministic regression test хийх нь хамгийн үр дүнтэй аргад тооцогдож байна.

\end{sloppy}
